\section{Perturbations of a retained exact inverse}
\label{sec:core-perturbations}

An inverse such as a real branch of the Lambert function can be useful as an exact building block even when it also has a convergent inverse--logarithmic expansion. Retaining that building block separates two different refinements: expansion of the core inverse itself, and correction for terms outside the core. The operation \wl{AsymptoticCoreInverse} implements the second refinement. Its truncation parameter counts powers of an auxiliary perturbation marker. The returned coefficients are complete functions of the exact core inverse; they are not individual terms of an exponent-sorted power--logarithmic series.

\subsection{The coefficient identity}

Choose the positive source coordinate $u\to0^+$ and write the original source point as
\[
 x=H(u),\qquad
 H(u)=x_0+\sigma u\quad\hbox{or}\quad H(u)=\frac{\sigma}{u},
 \qquad \sigma\in\{1,-1\}.
\]
Let $F_0(u)$ be the chosen core and $R(u)$ its perturbation in that coordinate. Let $u_0(y)$ be the selected exact inverse of $F_0$, and set $\phi(y)=H(u_0(y))$. Introduce a marker $\lambda$ by
\begin{equation}
 F_0(u_\lambda)+\lambda R(u_\lambda)=y,
 \qquad u_{\lambda=0}=u_0(y).
 \label{eq:core-marker-equation}
\end{equation}
At any fixed point with $F_0'(u_0)\ne0$, the inverse function theorem gives an analytic solution in $\lambda$ near zero whenever the data are analytic locally. The Lagrange--B\"urmann identity gives
\begin{equation}
 H(u_\lambda)=\phi(y)+\sum_{n\ge1}\lambda^n C_n(u_0(y)),
 \label{eq:core-marker-series}
\end{equation}
where
\begin{equation}
 C_n(u)=\frac{(-1)^n}{n!}
 \left(\frac{1}{F_0'(u)}\frac{d}{du}\right)^{n-1}
 \left(\frac{H'(u)R(u)^n}{F_0'(u)}\right).
 \label{eq:core-marker-coefficients}
\end{equation}
Indeed, with $t=F_0(u)$, the equation becomes
$t=y-\lambda R(u_0(t))$, and $d/dt=(F_0'(u))^{-1}d/du$.
Applying the ordinary Lagrange identity to the observable $\phi(t)$ proves
\eqref{eq:core-marker-coefficients}. This form of the identity differentiates
the explicit forward core. It avoids repeated symbolic differentiation of a
large expression containing \wl{ProductLog}.

For arbitrary data, this identity alone does not justify setting $\lambda=1$,
does not order the coefficients by their size as $u_0\to0$, and does not
bound the tail by its first omitted coefficient. The following restricted
class supplies those missing conclusions with a complete tail majorant.

\subsection{A class with an asymptotic tail bound}

Let $\M(u)=1+|\log u|$. Consider finite real power--logarithmic expressions
\begin{align}
 F_0(u)&=y_0+u^p Q_0(\log u)
       +\sum_{j=1}^{J}u^{p+\eta_j}Q_j(\log u),
       &p&\ne0,\quad \eta_j>0,\label{eq:core-admissible-core}\\
 R(u)&=\sum_{i=1}^{I}u^{p+\delta_i}P_i(\log u),
       &\delta_i&>0.\label{eq:core-admissible-perturbation}
\end{align}
All polynomials have real coefficients, $Q_0$ is nonzero, and all powers
are fixed real numbers. Equal powers are merged before the leading term
is selected. Put
\begin{equation}
 q=\deg Q_0,\qquad
 \delta=\min_i\delta_i,\qquad
 D=\max\left(0,\max_i\deg P_i-q\right),\qquad
 \chi(u)=u^\delta\M(u)^D.
 \label{eq:core-small-scale}
\end{equation}
For a zero perturbation the inverse remains exactly $\phi$, so these
parameters and the following tail estimate are unnecessary.

\begin{theorem}[An exact-core expansion with a complete marker tail]
\label{thm:core-marker-majorant}
For fixed data satisfying \eqref{eq:core-admissible-core}--\eqref{eq:core-admissible-perturbation}, the selected real inverse exists sufficiently near the source endpoint. Define $r_*=1$ for a finite source endpoint and $r_*=-1$ for an infinite source endpoint. There are positive constants $C,K$ and a source threshold such that, for every integer $N\ge0$ and all sufficiently small $u_0>0$ with $K\chi(u_0)<1$, the marker series converges at $\lambda=1$ and
\begin{equation}
 \left|F^{-1}(y)-\phi(y)-\sum_{n=1}^{N}C_n(u_0(y))\right|
 \le C u_0^{r_*}
 \frac{(K\chi(u_0))^{N+1}}{1-K\chi(u_0)}.
 \label{eq:core-geometric-tail}
\end{equation}
In particular, for every fixed $N$,
\begin{equation}
 F^{-1}(y)=\phi(y)+\sum_{n=1}^{N}C_n(u_0(y))
 +\OO\!\left(u_0^{r_*+(N+1)\delta}
                  \M(u_0)^{(N+1)D}\right).
 \label{eq:core-asymptotic-tail}
\end{equation}
The estimate concerns the complete omitted marker tail, including possible
cancellations in its first coefficient.
\end{theorem}
\begin{proof}
Write $u=u_0W$ and divide the equation
$F_0(u_0W)+\lambda R(u_0W)=F_0(u_0)$ by the leading scale
$u_0^p(\log u_0)^q$ and its nonzero leading coefficient.
Use $\tau=1/\log u_0$ and finitely many additional parameters for the
higher-power terms. For example,
\[
 \frac{(\log(u_0W))^k}{(\log u_0)^q}
 =(\log u_0)^{k-q}(1+\tau\log W)^k.
\]
If $k>q$, include the factor
$u_0^\eta(\log u_0)^{k-q}$ in its small parameter; if $k\le q$,
retain the nonnegative power $\tau^{q-k}$ as an analytic coefficient.
All such parameters tend to zero. Powers of $W$ and $\log W$ are analytic
in a fixed complex neighborhood of $W=1$ after fixing their branches.

At the limiting parameter origin, the normalized equation is $W^p-1=0$,
whose derivative at $W=1$ is $p\ne0$. The analytic implicit function theorem
therefore gives a holomorphic solution in a common polydisc. The core
parameters remain fixed when $\lambda$ varies, whereas every perturbation
parameter is multiplied by $\lambda$. Each perturbation parameter is
bounded in magnitude by a constant times $\chi(u_0)$ for small $u_0$.
Cauchy estimates in a smaller fixed polydisc consequently bound the total
coefficient of marker degree $n$ by $C K^n\chi(u_0)^n$. Equivalently,
the solution is holomorphic and uniformly bounded on a marker disc whose
radius is a fixed positive multiple of $1/\chi(u_0)$.

For a finite source endpoint, the nonconstant part of the observable is
$\sigma u_0W$; at infinity it is $\sigma u_0^{-1}W^{-1}$.
Both normalized observables are analytic near $W=1$, giving the additional
factor $u_0^{r_*}$. Summing the resulting geometric coefficient bound gives
\eqref{eq:core-geometric-tail}. The local real branch is unique, and
$F_0'(u)\sim p u^{p-1}Q_0(\log u)$ while $R'/F_0'\to0$,
so this analytic solution agrees with the selected real inverse near the
endpoint. Finally $\chi(u_0)\to0$, which proves
\eqref{eq:core-asymptotic-tail}.
\end{proof}

The constants in this theorem are existence bounds for fixed input data.
The implementation does not compute a numerical value of $C$, $K$, or the
source threshold. Its \wl{"MajorantContract"} therefore identifies the
bound as asymptotic existence and sets \wl{"NumericCertificate"} to
\wl{False}. The stated \wl{"SourceRadius"} controls a supplied inverse's
symbolic branch check; it is not asserted to be the unknown uniform
threshold in \eqref{eq:core-geometric-tail}.

The nonnegative value of $D$ in \eqref{eq:core-small-scale} is deliberately
conservative. A quotient by an extra power of $|\log u_0|$ can make an
individual coefficient smaller than the reported polynomial--log bound.
Keeping that improvement would require an additional logarithmic ordering
contract. Also, if the perturbation contains several different powers,
one complete marker coefficient may contain contributions lying beyond
the leading contribution of a later marker coefficient. Marker depth does
not promise a globally sorted list of all resulting monomials.

\subsection{The Lambert core and its first correction}

For the exact core
\[
 F_0(u)=y_0+a u^p(\log u+b)^q,\qquad p\ne0,
 \qquad q\in\Z_{\ge0},
\]
the parser verifies that the leading logarithmic polynomial is exactly the
displayed affine power. For $q=0$, the positive local inverse is
$u_0=((y-y_0)/a)^{1/p}$. For $q>0$, put
\[
 A=a(-1)^q,\qquad Y=\frac{y-y_0}{A}>0,\qquad
 k=-\frac pq,\qquad z=kY^{1/q}e^{pb/q}.
\]
Then
\begin{equation}
 u_0=Y^{1/p}\left(\frac{W_\nu(z)}{k}\right)^{-q/p},
 \qquad
 \nu=\begin{cases}-1,&p>0,\\0,&p<0.\end{cases}
 \label{eq:core-exact-lambert-inverse}
\end{equation}
These are the real branches for which $u_0\to0^+$. On the lower branch,
the real argument condition $-1/e\le z<0$ is recorded explicitly. The
quotient $W_\nu(z)/k$ is positive on the selected eventual branch, so all
displayed real powers have positive bases.

\begin{example}[Correcting $u\log u$ without expanding its exact inverse]
\label{ex:core-lambert-quadratic-perturbation}
For $F(u)=u\log u+u^2$ as $u\to0^+$, let
$u_0=y/W_{-1}(y)$. Formula~\eqref{eq:core-marker-coefficients} gives
\begin{align}
 C_1(u_0)&=-\frac{u_0^2}{1+\log u_0},\label{eq:core-lambert-first-correction}\\
 C_2(u_0)&=\frac{u_0^3(4\log u_0+3)}{2(1+\log u_0)^3}.
 \label{eq:core-lambert-second-correction}
\end{align}
Here $p=q=\delta=1$ and $D=0$. Thus truncation after $C_1$ has the
proved remainder $\OO(u_0^3)$, and truncation after $C_2$ has remainder
$\OO(u_0^4)$. The explicit omitted coefficient can exhibit a sharper
logarithmic scale, but the package reports the complete-tail bound
established above. Each displayed correction retains the exact
$W_{-1}$ expression in $u_0$.
\end{example}

\subsection{A declared forward remainder}

Suppose the actual forward function also contains an unknown error $E(u)$,
with the declared bounds
\begin{equation}
 E(u)=\OO(u^\rho\M(u)^k),\qquad
 E'(u)=\OO(u^{\rho-1}\M(u)^k),\qquad \rho>p.
 \label{eq:core-declared-input-error}
\end{equation}
The derivative bound is a hypothesis of this declaration, not a consequence
of the first bound alone. Since the retained forward derivative has scale
$u^{p-1}\M(u)^q$, the mean value argument of
Section~\ref{sec:transport} gives an additional error in the original
source variable of
\begin{equation}
 \OO\!\left(u_0^{r_*+\rho-p}
       \M(u_0)^{\max(0,k-q)}\right).
 \label{eq:core-input-error-transport}
\end{equation}
The gap $\rho>p$ ensures that this error is relatively small in the
positive source coordinate. Both the perturbed solution and the core
solution are asymptotic to the same $u_0$, so they can be used
interchangeably in these scale estimates. The output combines
\eqref{eq:core-input-error-transport} with the marker-tail bound by taking
the smaller source power, and the larger logarithmic degree if those
powers coincide. Increasing marker depth cannot improve this inherited
accuracy ceiling. The displayed marker coefficients still describe the
specified finite model; coefficients smaller than the unknown input error
are not asserted to be coefficients of every admissible actual function.

\subsection{The implemented operation and its evidence boundary}

\begin{lstlisting}[language=Mathematica]
AsymptoticCoreInverse[
  x Log[x], x^2, {x, 0}, {y, 2}]

AsymptoticCoreInverse[
  x + x^2, x^3, {x, 0}, {y, 1},
  "CoreInverse" -> (Sqrt[1 + 4 y] - 1)/2]
\end{lstlisting}

The first call automatically selects the lower real Lambert branch. The
second retains a supplied algebraic inverse after a bounded symbolic check
of $\phi(F_0(H(u)))=H(u)$ for $0<u<$\wl{"SourceRadius"} on the selected
side. An unverified inverse is rejected. Supplied parameter assumptions
must establish the reality of the coefficients and the eventual leading
sign. Exponents in this operation are exact real numeric constants; they
need not be rational.

The result identifies its scale as \wl{"ExactCorePerturbation"}.
Its kind is \wl{"CoreInverse"}, and \wl{"MarkerTerms"} stores pairs $\{n,C_n\}$,
with degree zero holding $\phi$. The full exact inverse of the core,
the local core inverse, the separately computed first omitted marker
coefficient, the branch certificate, and the tail contract are all retained.
The ordinary \wl{Normal} operation returns the retained expression. The
object does not claim an ordinary \wl{SeriesData} representation or an
ordinary exponent cutoff. Its independent regressions compare marker
coefficients with direct composition and exact algebraic inverses, check
both source orientations and infinite-endpoint error scaling, and compare
the Lambert correction with the original forward equation at high
precision. Such numerical checks supplement the theorem; they do not
turn its unspecified constants into a numerical certificate.

This implemented class requires strictly positive source-power gaps in
the perturbation. Perturbations of equal source power but smaller
logarithmic size, purely logarithmic leading cores, and general
exponential exact cores require additional scale and error-transport
contracts. They are not silently accepted by the marker identity.
The broader roadmap also includes refinement of the retained core,
composition of core objects with other expansion carriers, and computable
majorants. The present operation supplies the finite power--logarithmic
perturbation theorem and its exact-core implementation as one independently
testable stage.
